% Part: incompleteness
% Chapter: arithmetization-syntax
% Section: proofs-in-nd

\documentclass[../../../include/open-logic-section]{subfiles}

\begin{document}

\olfileid[pt]{inc}{art}{pnd}
\olsection{\usetoken{P}{derivation} em dedução natural}

\begin{explain}
Para aritmetizar !!{derivation}s, precisamos representar
!!{derivation}s como
números. Como !!{derivation}s são árvores de !!{formula}s nas quais
cada inferência possui um ou dois rótulos, uma representação recursiva
é a abordagem mais evidente: representamos !!a{derivation} como uma
tupla cujos componentes são o número de sub-!!{derivation}s imediatas
que levam às premissas da última inferência, as representações dessas
sub-!!{derivation}s, a !!{formula} final, o rótulo de descarga da última
inferência e um número que indica o tipo da última inferência.
\end{explain}

\begin{defn}
  Se $\delta$ é !!a{derivation} em dedução natural, então $\Gn{\delta}$
  é definido indutivamente da seguinte maneira:
  \begin{enumerate}
  \item 
    Se $\delta$ consiste apenas na hipótese~$!A$, então $\Gn{\delta}$ é
    $\tuple{0, \Gn{!A}, n}$. O número~$n$ é~$0$ se ela é uma hipótese
    !!{undischarged}; caso contrário, é o rótulo numérico.
  \item 
    Se $\delta$ termina em uma inferência com zero, uma, duas ou três
    premissas, então $\Gn{\delta}$ é
    \begin{align*}
      & \tuple{0, \Gn{!A}, n, k}, \\
      & \tuple{1, \Gn{\delta_1}, \Gn{!A}, n, k}, \\
      & \tuple{2, \Gn{\delta_1}, \Gn{\delta_2}, \Gn{!A}, n, k}, \text{ ou}\\
      & \tuple{3, \Gn{\delta_1}, \Gn{\delta_2}, \Gn{\delta_3}, \Gn{!A}, n,
        k},
    \end{align*}
    respectivamente. Aqui, $\delta_1$, $\delta_2$, $\delta_3$ são as
    sub-!!{derivation}s que terminam nas premissas da última inferência
    em $\delta$, $!A$ é a conclusão da última inferência em~$\delta$,
    $n$ é o rótulo de descarga da última inferência ($0$ se a inferência
    não descarrega nenhuma hipótese), e $k$~é dado pela tabela a seguir,
    conforme a regra usada na última inferência.
  
    \begin{tabular}{lccccccc}
      \text{Regra:} & \Intro{\land} & \Elim{\land} & \Intro{\lor} & \Elim{\lor} \\
      $k$: & 1 & 2 & 3 & 4  \\[1.5ex]
      \text{Regra:} &  \Intro{\lif} & \Elim{\lif} & \Intro{\lnot} & \Elim{\lnot} \\
      $k$: & 5 & 6 & 7 & 8 \\[1.5ex]
      \text{Regra:}  &  \FalseInt & \FalseCl & \Intro{\lforall} & \Elim{\lforall} \\
      $k$: & 9 & 10 & 11 & 12 \\[1.5ex]
      \text{Regra:} & \Intro{\lexists} & \Elim{\lexists} & \Intro{\eq} & \Elim{\eq} \\
      $k$: & 13 & 14 & 15 & 16 
    \end{tabular}
  \end{enumerate}
\end{defn}

\begin{ex}
  Considere a !!{derivation} muito simples
  \begin{prooftree}
    \AxiomC{$\Discharge{!A \land !B}{1}$}
    \RightLabel{\Elim{\land}}
    \UnaryInfC{$!A$}
    \DischargeRule{\Intro{\lif}}{1}
    \UnaryInfC{$(!A \land !B) \lif !A$}
  \end{prooftree}
  O número de G\"odel da hipótese seria $d_0 = \tuple{0,
    \Gn{!A \land !B}, 1}$. O número de G\"odel da !!{derivation} que
  termina na conclusão de~$\Elim{\land}$ seria $d_1 = \tuple{1,
     d_0, \Gn{!A}, 0, 2}$ ($1$ porque $\Elim{\land}$ tem uma premissa,
  o número de G\"odel da conclusão~$!A$, $0$ porque nenhuma hipótese é
  descarregada, e $2$ é o número que codifica $\Elim{\land}$). O número
  de G\"odel da !!{derivation} inteira é, então,
  $\tuple{1, d_1, \Gn{((!A \land !B) \lif
      !A)}, 1, 5}$, isto é,
  \[
  \tuple{1, \tuple{1, \tuple{0, \Gn{(!A \land !B)}, 1}, \Gn{!A}, 0, 2},
    \Gn{((!A \land !B) \lif !A)}, 1, 5}.
  \]
\end{ex}

\begin{explain}
Escolhida uma representação das !!{derivation}s, precisamos também
mostrar que podemos manipular recursivamente de modo primitivo os
números de G\"odel dessas !!{derivation}s e expressar suas propriedades
e relações essenciais. Algumas operações são simples: por exemplo,
dado um número de G\"odel~$d$ de !!a{derivation},
$\fn{EndFmla}(d) = (d)_{(d)_0+1}$ nos dá o número de G\"odel de sua
!!{formula} final, $\fn{DischargeLabel}(d) = (d)_{(d)_0 + 2}$ nos dá o
rótulo de descarga, e $\fn{LastRule}(d) = (d)_{(d)_0 + 3}$, o número
que indica o tipo da última inferência. Algumas propriedades são muito
mais difíceis. Esboçaremos ao menos como tratá-las. O objetivo é mostrar
que a relação ``$\delta$ é !!a{derivation} de~$!A$ a partir
de~$\Gamma$'' é uma relação recursiva primitiva dos números de G\"odel
de $\delta$ e~$!A$.
\end{explain}

\begin{prop}
  As relações a seguir são recursivas primitivas:
  \begin{enumerate}
  \item $!A$ ocorre como hipótese em~$\delta$ com rótulo~$n$.
  \item Todas as hipóteses em $\delta$ com rótulo~$n$ têm a forma~$!A$
    (isto é, podemos !!{discharge} a hipótese~$!A$ usando o rótulo~$n$
    em~$\delta$).
  \end{enumerate}
\end{prop}

\begin{proof}
  Precisamos mostrar que as relações correspondentes entre números de
  G\"odel de !!{formula}s e números de G\"odel de !!{derivation}s são
  recursivas primitivas.
  \begin{enumerate}
  \item Queremos mostrar que $\fn{Assum}(x, d, n)$, que vale se $x$ é
    o número de G\"odel de uma hipótese da !!{derivation} com número de
    G\"odel~$d$ rotulada por~$n$, é recursiva primitiva. Isso ocorre se
    a !!{derivation} com número de G\"odel~$\tuple{0, x, n}$ é uma
    sub-!!{derivation} de~$d$. Observe que a maneira como codificamos
    !!{derivation}s é um caso especial da codificação de árvores
    introduzida em \olref[cmp][rec][tre]{sec}; portanto, a função
    recursiva primitiva $\fn{SubtreeSeq}(d)$ fornece uma sequência de
    números de G\"odel de todas as sub-!!{derivation}s de~$d$ (de
    comprimento no máximo $d$). Assim, podemos definir
    \[
    \fn{Assum}(x, d, n) \defiff \bexists{i<d}{(\fn{SubtreeSeq}(d))_i =
      \tuple{0, x, n}}.
    \]
  \item Queremos mostrar que $\fn{Discharge}(x, d, n)$, que vale se
    todas as hipóteses com rótulo~$n$ na !!{derivation} com número de
    G\"odel~$d$ são a !!{formula} com número de G\"odel~$x$. Mas essa
    relação vale sse $\bforall{y<d}{(\fn{Assum}(y, d, n) \lif y = x)}$.
  \end{enumerate}
\end{proof}

\begin{prop}
  \ollabel{prop:followsby}
  A propriedade $\fn{Correct}(d)$, que vale sse a última inferência na
  !!{derivation}~$\delta$ com número de G\"odel~$d$ é correta, é
  recursiva primitiva.
\end{prop}

\begin{proof}
  Aqui precisamos mostrar que, para cada regra de inferência~$R$, a
  relação $\fn{FollowsBy}_R(d)$ é recursiva primitiva, em que
  $\fn{FollowsBy}_R(d)$ vale sse $d$ é o número de G\"odel de
  !!a{derivation}~$\delta$ e a !!{formula} final de~$\delta$ decorre,
  por uma aplicação correta de~$R$, das sub-!!{derivation}s imediatas
  de~$\delta$.

  Um caso simples é o da regra \Intro{\land}. Se $\delta$ termina em
  uma inferência $\Intro{\land}$ correta, ela tem esta forma:
  \begin{prooftree}
    \AxiomC{}
    \RightLabel{$\delta_1$}
    \DeduceC{$!A$}
    
    \AxiomC{}
    \RightLabel{$\delta_2$}
    \DeduceC{$!B$}

    \RightLabel{\Intro\land}
    \BinaryInfC{$!A \land !B$}
  \end{prooftree}
  Então o número de G\"odel~$d$ de~$\delta$ é $\tuple{2, d_1, d_2,
    \Gn{(!A \land !B)}, 0, k}$, em que
  $\fn{EndFmla}(d_1) = \Gn{!A}$, $\fn{EndFmla}(d_2) = \Gn{!B}$,
  $n=0$ e $k=1$. Assim, podemos definir
  $\fn{FollowsBy}_{\Intro\land}(d)$ como
  \begin{multline*}
    (d)_0 = 2 \land \fn{DischargeLabel}(d) = 0 \land \fn{LastRule}(d) = 1 \land {}\\
    \fn{EndFmla}(d) = {}\\
    \Gn{(} \concat \fn{EndFmla}((d)_1) \concat \Gn{\land}
    \concat \fn{EndFmla}((d)_2) \concat \Gn{)}.
  \end{multline*}

  Outro exemplo simples é a regra $\Intro\eq$. Ela não tem premissas,
  de modo que $(d)_0 = 0$, como as hipóteses. Também não tem rótulo de
  descarga, isto é, $n=0$. Contudo, $!A$ deve ter a forma $\eq[t][t]$
  para um termo fechado~$t$. Aqui, uma definição recursiva primitiva é
  \begin{multline*}
    (d)_0 = 0 \land \fn{DischargeLabel}(d) = 0 \land {}\\
    \bexists{t<d}{(\fn{ClTerm}(t) \land 
      \fn{EndFmla}(d) = {}\\
      \Gn{{\eq}(} \concat t \concat \Gn{,} \concat t \concat \Gn{)})}.
  \end{multline*}

  Como exemplo mais complicado, $\fn{FollowsBy}_{\Intro{\lif}}(d)$
  vale sse a !!{formula} final de~$\delta$ tem a forma $(!A \lif !B)$,
  em que a !!{formula} final de $\delta_1$ é~$!B$ e toda hipótese
  em~$\delta$ rotulada por~$n$ tem a forma~$!A$. Podemos expressar isso
  recursivamente de modo primitivo por
  \begin{multline*}
    (d)_0 = 1 \land {}\\
    \bexists{a<d}{(\fn{Discharge}(a, (d)_1, \fn{DischargeLabel}(d)) \land {}}\\
      \fn{EndFmla}(d) = (\Gn{(} \concat a \concat \Gn{\lif}
      \concat \fn{EndFmla}((d)_1) \concat \Gn{)}))
  \end{multline*}
  (Pense em $a$ como o número de G\"odel de~$!A$.)

  Como outro exemplo, considere \Intro{\lexists}. Aqui, a última
  inferência em~$\delta$ é correta sse existem !!a{formula}~$!A$, um
  termo fechado~$t$ e !!a{variable}~$x$ tais que $\Subst{!A}{t}{x}$ é
  a !!{formula} final da !!{derivation}~$\delta_1$ e
  $\lexists[x][!A]$ é a conclusão da última inferência. Logo,
  $\fn{FollowsBy}_{\Intro{\lexists}}(d)$ vale sse
  \begin{multline*}
    (d)_0 = 1 \land \fn{DischargeLabel}(d) = 0 \land {} \\
    \bexists{a < d}{\bexists{x<d}{\bexists{t<d}{
          (\fn{ClTerm}(t) \land \fn{Var}(x)  \land {}}}}\\
    \fn{Subst}(a,t,x) = \fn{EndFmla}((d)_1) \land
    \fn{EndFmla}(d) = (\Gn{\lexists} \concat x \concat a)).
  \end{multline*}

  Definimos então $\fn{Correct}(d)$ como
  \begin{multline*}
    \fn{Sent}(\fn{EndFmla}(d)) \land {}\\
    (\fn{LastRule}(d) = 1 \land
    \fn{FollowsBy}_{\Intro\land}(d)) \lor \dots \lor {}\\
    (\fn{LastRule}(d) = 16 \land \fn{FollowsBy}_{\Elim\eq}(d)) \lor {}\\
    \bexists{n<d}{\bexists{x<d}{(d = \tuple{0, x, n})}}.
  \end{multline*}
  A primeira linha garante que a !!{formula} final de~$d$ é uma
  sentença. A última linha cobre o caso em que $d$ é apenas uma
  hipótese.
\end{proof}

\begin{prob}
  Defina as propriedades a seguir como em
  \olref[inc][art][pnd]{prop:followsby}:
  \begin{enumerate}
  \item $\fn{FollowsBy}_{\Elim{\lif}}(d)$,
  \item $\fn{FollowsBy}_{\Elim{\eq}}(d)$,
  \item $\fn{FollowsBy}_{\Elim{\lor}}(d)$,
  \item $\fn{FollowsBy}_{\Intro{\lforall}}(d)$.
  \end{enumerate}
  Para a última, você também terá de mostrar que é possível testar
  recursivamente de modo primitivo se a última inferência da
  !!{derivation} com número de G\"odel~$d$ satisfaz a condição da
  variável própria, isto é, se a variável própria~$a$ da inferência
  $\Intro{\lforall}$ não ocorre nem na !!{formula} final de~$d$ nem em
  uma hipótese aberta de~$d$. Você pode usar o predicado recursivo
  primitivo $\fn{OpenAssum}$ de
  \olref[inc][art][pnd]{prop:openassum} para isso.
\end{prob}

\begin{prop}
  \ollabel{prop:deriv}
  A relação $\fn{Deriv}(d)$, que vale se $d$ é o número de G\"odel de
  !!a{derivation} correta~$\delta$, é recursiva primitiva.
\end{prop}

\begin{proof}
  !!^a{derivation}~$\delta$ é correta se cada uma de suas inferências
  é uma aplicação correta de uma regra, isto é, se cada uma de suas
  sub-!!{derivation}s termina em uma inferência correta. Assim,
  $\fn{Deriv}(d)$ sse
  \[
  \bforall{i<\len{\fn{SubtreeSeq}(d)}}{\fn{Correct}((\fn{SubtreeSeq}(d))_i)}
  \]
\end{proof}

\begin{prop}
  \ollabel{prop:openassum} A relação $\fn{OpenAssum}(z, d)$, que vale
  se $z$ é o número de G\"odel de !!a{undischarged} hipótese~$!A$ da
  !!{derivation}~$\delta$ com número de G\"odel~$d$, é recursiva
  primitiva.
\end{prop}

\begin{proof}
  Uma ocorrência de uma hipótese é !!{discharged} se ocorre com
  rótulo~$n$ em uma sub-!!{derivation} de~$\delta$ que termina em uma
  regra com rótulo de descarga~$n$. Assim, $!A$ é !!a{undischarged}
  hipótese de~$\delta$ se ao menos uma de suas ocorrências não é
  !!{discharged} em~$\delta$. Precisamos ter cuidado: $\delta$ pode
  conter ocorrências de~$!A$ tanto !!{discharged} quanto
  !!{undischarged}.

  Considere uma sequência $\delta_0$, \dots, $\delta_k$ na qual
  $\delta_0 = \delta$, $\delta_k$ é a hipótese $\Discharge{!A}{n}$
  (para algum~$n$) e $\delta_{i+1}$ é uma sub-!!{derivation} imediata
  de~$\delta_i$. Se existe uma sequência assim na qual nenhum
  $\delta_i$ termina em uma inferência com rótulo de descarga~$n$,
  então $!A$ é !!a{undischarged} hipótese de~$\delta$.

  A função recursiva primitiva $\fn{SubtreeSeq}(d)$ nos fornece uma
  sequência de números de G\"odel de todas as sub-!!{derivation}s
  de~$\delta$. Toda sequência de números de G\"odel de
  sub-!!{derivation}s de~$\delta$ é uma subsequência dela. Ser uma
  subsequência é uma relação recursiva primitiva:
  $\fn{Subseq}(s, s')$ vale sse
  $\bforall{i<\len{s}}{\lexists[j<\len{s'}][(s)_i = (s')_j]}$. Ser uma
  sub-!!{derivation} imediata também é: $\fn{Subderiv}(d, d')$ sse
  $\bexists{j<(d')_0}{d = (d')_j}$. Assim, podemos definir
  $\fn{OpenAssum}(z, d)$ por
  \begin{multline*}
    \bexists{s<\fn{SubtreeSeq}(d)}{(\fn{Subseq}(s, \fn{SubtreeSeq}(d))
      \land (s)_0 = d \land {}} \\
    \bexists{n<d}{((s)_{\len{s} \tsub 1} = \tuple{0, z, n} \land {}}\\
      \bforall{i<(\len{s} \tsub 1)}{(\fn{Subderiv}((s)_{i+1}, (s)_i) \land {}}\\
      \fn{DischargeLabel}((s)_i) \neq n))).
  \end{multline*}
\end{proof}

\begin{prop}
  \ollabel{prop:prf-prim-rec}
  Suponha que $\Gamma$ seja um conjunto recursivo primitivo de
  !!{sentence}s. Então a relação $\Prf[\Gamma](x, y)$ que expressa
  ``$x$ é o código de !!a{derivation}~$\delta$ de $!A$ a partir de
  hipóteses !!{undischarged} em~$\Gamma$, e $y$ é o número de G\"odel
  de~$!A$'' é recursiva primitiva.
\end{prop}

\begin{proof}
  Suponha que ``$y \in \Gamma$'' seja dada pelo predicado recursivo
  primitivo~$R_\Gamma(y)$. Precisamos mostrar que
  $\Prf[\Gamma](x, y)$, que vale sse $y$ é o número de G\"odel de uma
  sentença~$!A$ e $x$~é o código de uma !!{derivation} em dedução
  natural com !!{formula} final~$!A$ e todas as hipóteses
  !!{undischarged} em~$\Gamma$, é recursiva primitiva.

  Por \olref{prop:deriv}, a propriedade $\fn{Deriv}(x)$, que vale sse
  $x$ é o número de G\"odel de !!a{derivation} correta~$\delta$ em
  dedução natural, é recursiva primitiva. Assim, podemos definir
  $\Prf[\Gamma](x, y)$ por
  \begin{align*}
    \Prf[\Gamma](x, y) \defiff {}
    & \fn{Deriv}(x) \land \fn{EndFmla}(x) = y \land {} \\
    & \bforall{z < x}{(\fn{OpenAssum}(z, x) \lif R_\Gamma(z))}.
  \end{align*}
\end{proof}

\end{document}
